% Storyline 1, paragraph 4 (round 2): approach, opening with the positive concept.
Here, we develop a model of indirect reciprocity in which reputations have inertia, and in which this inertia can differ between losing and regaining standing (Fig.~1). An established reputation changes only after a run of contrary assessments. Negative assessments move an individual in good standing through a sequence of warning states before the reputation turns bad, and positive assessments move an individual in bad standing through recovery states before it turns good. The two paths can differ in length, so that losing and regaining standing may require different amounts of evidence. Each intermediate state is interpreted as good or bad with a probability that sets how strongly incomplete evidence already affects treatment. The model thus treats two properties of reputational change as separate design choices, namely how much evidence is needed to change a reputation in each direction and how incomplete evidence is expressed in the meantime.
